% GENERATED FILE. Edit canonical lessons, then run npm run generate:books.
% canonical-source-hash: fnv1a64:357fe31375a781d8
% canonical-lessons: GU-C204-samaypatrak, GU-C204-puchparach, GU-C204-pletphorm, GU-C204-upadvu, GU-C204-chutvu

\chapter{Reading the timetable}
\label{ch:gu-reading-the-timetable}
\hlchaptermodality{204}

\begin{chapteropening}
\textbf{By the end of this chapter:} \emph{I can read the timetable and say when the train left.}

Complete the last lesson of chapter 204: I can read the timetable and say when the train left.
\end{chapteropening}

\section[\texorpdfstring{\gu{સમયપત્રક} (samaypatrak)}{સમયપત્રક (samaypatrak)}]{\gu{સમયપત્રક} (samaypatrak) — timetable}

\label{lesson:GU-C204-samaypatrak}



\begin{quote}
Before the new one: say the Gujarati for to drop off, then the Gujarati for to see off.
\end{quote}

\subsection*{You'll want to know: \gu{સમયપત્રક}}
\textbf{\gu{સમયપત્રક}} — \emph{samaypatrak} — \textquotedblleft{}timetable\textquotedblright{}.

\begin{quote}
\textbf{\gu{સ્ટેશન} \gu{પર} \gu{નવું} \gu{સમયપત્રક} \gu{હતું}.} — \emph{sṭeśan par navũ samaypatrak hatũ} — \textquotedblleft{}There was a new timetable at the station.\textquotedblright{}
\end{quote}

Built on \textbf{\gu{સમય}}, \emph{samay}, 'time'.

\begin{cousinweb}
\textbf{\gu{સમયપત્રક}} (\emph{samaypatrak}) is a learned compound built from Sanskrit \emph{samaya}, \textquotedblleft{}time\textquotedblright{}, and \emph{patraka}, \textquotedblleft{}leaf, sheet\textquotedblright{}, a form of \emph{páttra}, \textquotedblleft{}leaf\textquotedblright{} (CDIAL 7733).
\end{cousinweb}

\begin{grammarlens}[title={What you've built}]
One more word for everyday life at A2.
\end{grammarlens}

\subsection*{Your turn}
\begin{itemize}
\raggedright
  \item \emph{Say it:} \emph{samaypatrak}
  \item \emph{Say it:} \emph{samaypatrak}, once more
  \item \emph{Say it:} \emph{sṭeśan par navũ samaypatrak hatũ} — There was a new timetable at the station.
  \item \emph{Recall:} say the Gujarati for to drop off, then the Gujarati for to see off, then say \emph{samaypatrak} again
\end{itemize}

\begin{tcolorbox}[breakable,colback=teal!4,colframe=teal!35!black,title={Before you move on}]
What does \gu{સમયપત્રક} mean? (\textquotedblleft{}Timetable\textquotedblright{}.) Say it once more.
\end{tcolorbox}
\section[\texorpdfstring{\gu{પૂછપરછ} (pūchparach)}{પૂછપરછ (pūchparach)}]{\gu{પૂછપરછ} (pūchparach) — enquiry, enquiries}

\label{lesson:GU-C204-puchparach}



\begin{quote}
Before the new one: say the Gujarati for to see off, then the Gujarati for timetable.
\end{quote}

\subsection*{You'll want to know: \gu{પૂછપરછ}}
\textbf{\gu{પૂછપરછ}} — \emph{pūchparach} — \textquotedblleft{}enquiry, enquiries\textquotedblright{}.

\begin{quote}
\textbf{\gu{મેં} \gu{પૂછપરછમાં} \gu{સમય} \gu{પૂછ્યો}.} — \emph{mẽ pūchparachmā̃ samay pūchyo} — \textquotedblleft{}I asked the time at the enquiry desk.\textquotedblright{}
\end{quote}

A feminine noun.

\begin{cousinweb}
\textbf{\gu{પૂછપરછ}} (\emph{pūchparach}) is built on \emph{pūch}, \textquotedblleft{}asking\textquotedblright{}, doubled with a rhyming partner. CDIAL lists Gujarati \emph{puch-parach}, \textquotedblleft{}questioning\textquotedblright{}, under Sanskrit \emph{pṛcchā}, \textquotedblleft{}asking\textquotedblright{} (CDIAL 8354), the noun beside \emph{pṛccháti}, \textquotedblleft{}asks\textquotedblright{}, which gave \emph{pūchvũ}, \textquotedblleft{}to ask\textquotedblright{} (CDIAL 8352).
\end{cousinweb}

\begin{grammarlens}[title={What you've built}]
One more word for everyday life at A2.
\end{grammarlens}

\subsection*{Your turn}
\begin{itemize}
\raggedright
  \item \emph{Say it:} \emph{pūchparach}
  \item \emph{Say it:} \emph{pūchparach}, once more
  \item \emph{Say it:} \emph{mẽ pūchparachmā̃ samay pūchyo} — I asked the time at the enquiry desk.
  \item \emph{Recall:} say the Gujarati for to see off, then the Gujarati for timetable, then say \emph{pūchparach} again
\end{itemize}

\begin{tcolorbox}[breakable,colback=teal!4,colframe=teal!35!black,title={Before you move on}]
What does \gu{પૂછપરછ} mean? (\textquotedblleft{}Enquiry, enquiries\textquotedblright{}.) Say it once more.
\end{tcolorbox}
\section[\texorpdfstring{\gu{પ્લેટફોર્મ} (pleṭphorm)}{પ્લેટફોર્મ (pleṭphorm)}]{\gu{પ્લેટફોર્મ} (pleṭphorm) — platform}

\label{lesson:GU-C204-pletphorm}



\begin{quote}
Before the new one: say the Gujarati for timetable, then the Gujarati for enquiry.
\end{quote}

\subsection*{You'll want to know: \gu{પ્લેટફોર્મ}}
\textbf{\gu{પ્લેટફોર્મ}} — \emph{pleṭphorm} — \textquotedblleft{}platform\textquotedblright{}.

\begin{quote}
\textbf{\gu{ટ્રેન} \gu{ત્રીજા} \gu{પ્લેટફોર્મ} \gu{પર} \gu{આવી}.} — \emph{ṭren trījā pleṭphorm par āvī} — \textquotedblleft{}The train came in at platform three.\textquotedblright{}
\end{quote}

\begin{cousinweb}
\textbf{\gu{પ્લેટફોર્મ}} (\emph{pleṭphorm}) is a loan from English \emph{platform}.
\end{cousinweb}

\begin{grammarlens}[title={What you've built}]
One more word for everyday life at A2.
\end{grammarlens}

\subsection*{Your turn}
\begin{itemize}
\raggedright
  \item \emph{Say it:} \emph{pleṭphorm}
  \item \emph{Say it:} \emph{pleṭphorm}, once more
  \item \emph{Say it:} \emph{ṭren trījā pleṭphorm par āvī} — The train came in at platform three.
  \item \emph{Recall:} say the Gujarati for timetable, then the Gujarati for enquiry, then say \emph{pleṭphorm} again
\end{itemize}

\begin{tcolorbox}[breakable,colback=teal!4,colframe=teal!35!black,title={Before you move on}]
What does \gu{પ્લેટફોર્મ} mean? (\textquotedblleft{}Platform\textquotedblright{}.) Say it once more.
\end{tcolorbox}
\section[\texorpdfstring{\gu{ઊપડવું} (ūpaḍvũ)}{ઊપડવું (ūpaḍvũ)}]{\gu{ઊપડવું} (ūpaḍvũ) — to start off, to depart}

\label{lesson:GU-C204-upadvu}



\begin{quote}
Before the new one: say the Gujarati for enquiry, then the Gujarati for platform.
\end{quote}

\subsection*{You'll want to know: \gu{ઊપડવું}}
\textbf{\gu{ઊપડવું}} — \emph{ūpaḍvũ} — \textquotedblleft{}to start off, to depart\textquotedblright{}.

\begin{quote}
\textbf{\gu{ટ્રેન} \gu{સવારે} \gu{ઊપડી}.} — \emph{ṭren savāre ūpaḍī} — \textquotedblleft{}The train left in the morning.\textquotedblright{}
\end{quote}

\begin{cousinweb}
\textbf{\gu{ઊપડવું}} (\emph{ūpaḍvũ}) is a tadbhava. CDIAL lists it as \emph{upaṛvũ}, \textquotedblleft{}to move away, depart, be lifted, rise\textquotedblright{}, under Sanskrit \emph{útpatati}, \textquotedblleft{}flies up, ascends\textquotedblright{} (CDIAL 1810).
\end{cousinweb}

\begin{grammarlens}[title={What you've built}]
One more word for everyday life at A2.
\end{grammarlens}

\subsection*{Your turn}
\begin{itemize}
\raggedright
  \item \emph{Say it:} \emph{ūpaḍvũ}
  \item \emph{Say it:} \emph{ūpaḍvũ}, once more
  \item \emph{Say it:} \emph{ṭren savāre ūpaḍī} — The train left in the morning.
  \item \emph{Recall:} say the Gujarati for enquiry, then the Gujarati for platform, then say \emph{ūpaḍvũ} again
\end{itemize}

\begin{tcolorbox}[breakable,colback=teal!4,colframe=teal!35!black,title={Before you move on}]
What does \gu{ઊપડવું} mean? (\textquotedblleft{}To start off, to depart\textquotedblright{}.) Say it once more.
\end{tcolorbox}
\section[\texorpdfstring{\gu{છૂટવું} (chūṭvũ)}{છૂટવું (chūṭvũ)}]{\gu{છૂટવું} (chūṭvũ) — to leave, to get loose}

\label{lesson:GU-C204-chutvu}



\begin{quote}
Before the new one: say the Gujarati for platform, then the Gujarati for to start off.
\end{quote}

\subsection*{You'll want to know: \gu{છૂટવું}}
\textbf{\gu{છૂટવું}} — \emph{chūṭvũ} — \textquotedblleft{}to leave, to get loose\textquotedblright{}.

\begin{quote}
\textbf{\gu{મારી} \gu{ટ્રેન} \gu{છૂટી} \gu{ગઈ}.} — \emph{mārī ṭren chūṭī gaī} — \textquotedblleft{}My train left without me.\textquotedblright{}
\end{quote}

\begin{cousinweb}
\textbf{\gu{છૂટવું}} (\emph{chūṭvũ}) is a tadbhava. CDIAL lists it, glossed \textquotedblleft{}to get loose\textquotedblright{}, under \emph{*kṣuṭyatē}, \textquotedblleft{}is released\textquotedblright{}, a reconstructed form (CDIAL 3707).
\end{cousinweb}

\begin{grammarlens}[title={What you've built}]
One more word for everyday life at A2.
\end{grammarlens}

\subsection*{Your turn}
\begin{itemize}
\raggedright
  \item \emph{Say it:} \emph{chūṭvũ}
  \item \emph{Say it:} \emph{chūṭvũ}, once more
  \item \emph{Say it:} \emph{mārī ṭren chūṭī gaī} — My train left without me.
  \item \emph{Recall:} say the Gujarati for platform, then the Gujarati for to start off, then say \emph{chūṭvũ} again
\end{itemize}

\begin{tcolorbox}[breakable,colback=teal!4,colframe=teal!35!black,title={Before you move on}]
What does \gu{છૂટવું} mean? (\textquotedblleft{}To leave, to get loose\textquotedblright{}.) Say it once more.
\end{tcolorbox}
